\subsubsection{Axio-Dilaton Dark Energy: A Dynamical Systems and Bayesian Inference Analysis --- arXiv:2609.22539}

\textbf{Authors:} Mario Ramos-Hamud, Gabriela García-Arroyo, Fernando Quevedo, J. Alberto Vázquez

\paragraph{主な主張}
曲がった field-space metric を持つ axion--dilaton 系について、一般化 Albrecht--Skordis potential の力学と SNe Ia・BAO・Planck 2018 distance-prior による制限を解析し、steep な potential ほど late-time acceleration を維持するために強い kinetic coupling を要する縮退が現れ、現データでは dark energy を記述できるものの $\Lambda$CDM に対する統計的優位は小さいことを示した \cite{RamosHamud:2026axiodilaton}。

\paragraph{新規性}
supergravity・string compactification で動機づけられる axion--dilaton の curved field space と、漸近指数 potential からの polynomial correction を同時に取り込み、dynamical-system 解析・数値発展・Bayesian inference を一つの枠組みで接続した点にある。

\paragraph{位置付け}
hyperbolic/curved field-space を持つ modulus cosmology で、potential の steepness と kinetic coupling の競合が late-time dark energy をどう実現し観測制約を受けるかを調べる benchmark であり、axion--saxion tracker 系の研究を現在宇宙の加速膨張へ接続する位置にある。
